\documentclass[conference]{IEEEtran}
\usepackage[T1]{fontenc}
\usepackage{cite}
\usepackage{amsmath}
\usepackage{textcomp}

\begin{document}

\title{Automated Microscopic Diagnosis of Malaria}

\author{\IEEEauthorblockN{Agrapujya Lashkari}
\and
\IEEEauthorblockN{Deepanshu Mehta}}

\maketitle

% ---------------------------------------------------------------------------
% Section I was rewritten in Milestone 3. Sections II-VIII are the Phase-1
% literature survey, unchanged.
% ---------------------------------------------------------------------------
\section{Introduction}
Malaria remains one of the largest infectious-disease burdens in the world, with an
estimated 249 million cases and 608{,}000 deaths in 2022 alone~\cite{who2023}. Light
microscopy of Giemsa-stained blood smears is still the reference diagnostic in the
low-resource settings where most of those cases occur.
Analyzing malaria through manual microscopy is intensely time-consuming. It also
introduces significant inter-observer variability, which naturally drove the research
community to explore automated computer vision alternatives. When the NIH released its
Malaria Dataset---containing over 27,000 images of individual cells---it finally gave
developers the foundational raw material needed to start training diagnostic deep
learning models.

The work that followed (Sections~\ref{sec:transfer}--\ref{sec:danet}) moved in
two directions that pull against each other. One line chased accuracy, from transfer
learning and ensembles~\cite{rajaraman2018,rajaraman2019} to capsule networks and
compact transformers~\cite{madhu2023,masud2022}. The other chased deployability, with
small custom CNNs for phones and edge devices~\cite{rahman2019,fuhad2020,masud2020}.
Reported cell-level accuracies now sit between 95\% and 99\%. Most of these numbers,
however, come from one evaluation habit: the NIH cells are split \emph{at random},
so cells from the same patient appear in both the training and the test set
(Section~\ref{sec:gaps}). A score obtained this way measures how well a model
recognizes patients it has already seen, not how it would perform on a new patient's
slide.

In this project we take DANet~\cite{ali2025danet} (Section~\ref{sec:danet}), a 2.3~M-parameter dilated attention
network that reports 97.95\% accuracy on NIH with public code, as a representative of
the lightweight line. We first reproduce it under the paper's own protocol.
Reproduction required repairing the released code, which fails on import and has a
32-way output layer for a 2-class task. It also required resolving an optimizer
inconsistency: the learning rate given in the paper's hyperparameter table does not
train at all. Our reproduction reaches 97.02\% accuracy, 96.30\% sensitivity and
0.9945 ROC-AUC. That is inside the paper's own 5-fold cross-validation range
(96.97--97.42\%) but below its headline figure, and the shortfall is concentrated in
sensitivity, i.e.\ in missed infections.

Analyzing the reproduced baseline exposes the mechanism behind the evaluation concern
above. Under the random split, all 201 test patients also contribute training cells,
and the patient identity alone explains 74--90\% of the variance in a cell's mean
color. Each slide's stain therefore acts as a near-unique fingerprint that a network
can exploit instead of parasite morphology. Building on evidence from computational
pathology that color augmentation improves robustness to stain
variation~\cite{tellez2019}, we hypothesize that \emph{adding stain-color
augmentation, calibrated to the measured between-patient color spread, will raise
DANet's sensitivity on unseen patients}. Specifically, we predict a gain of $+0.5$ to
$+1.5$ points in sensitivity under a patient-level split, and no gain ($\pm0.3$ points)
under the leaky cell-level split. We fix these predictions before running the
experiment.

The contributions of this work are:
\begin{itemize}
\item a documented reproduction of DANet, including the code and optimizer fixes
needed for it to train and an analysis of the reproduction gap;
\item a quantification of patient leakage and stain fingerprinting in the NIH dataset
under the commonly used cell-level protocol;
\item a single-variable, pre-registered hypothesis with predicted effect sizes, tested
on a patient-level split with sensitivity, balanced accuracy and ROC-AUC as the
primary metrics. No clinical claims are made.
\end{itemize}
The remainder of this report reviews prior work by theme
(Sections~\ref{sec:transfer}--\ref{sec:danet}) and summarizes the open research
gaps that motivate our hypothesis (Section~\ref{sec:gaps}).

% ======================= Phase-1 survey (unchanged) =========================
\section{Transfer Learning and Ensemble Approaches}\label{sec:transfer}
In the early days of automated malaria detection, nobody really had the massive
datasets required to train deep neural networks from the ground up without running
into severe overfitting. To navigate this limitation, Rajaraman \emph{et al.} (2018)
turned to transfer learning. They took pre-trained convolutional neural networks
(CNNs), such as ResNet-50, and repurposed them to handle complex feature
extraction~\cite{rajaraman2018}.

\textbf{Problem Solved:} The sheer inability to train deep models from scratch due to
data scarcity, alongside a desire to avoid tedious manual feature engineering.\\
\textbf{How it was Solved:} The researchers brought in pre-trained ImageNet CNNs to
act as automatic feature extractors.\\
\textbf{Limitations:} Single deep learning models exhibited severe prediction
variance. Basically, they proved far too sensitive to minor fluctuations and quirks
within the training set.

To smooth out this instability, Rajaraman \emph{et al.} (2019) proposed a model
averaging ensemble~\cite{rajaraman2019}.

\textbf{Problem Solved:} The jumpy, unreliable nature of relying on a single
transfer-learning model for medical diagnosis.\\
\textbf{How it was Solved:} By blending fine-tuned VGG-19 and SqueezeNet
architectures. Because these distinct networks tended to make mistakes in completely
different areas, averaging their output probabilities effectively stabilized the whole
system.\\
\textbf{Limitations:} Combining models is computationally expensive. The resulting
setup was far too heavy for low-power clinics in the developing world where
point-of-care malaria testing actually happens.

\section{Resource-Constrained Architectures}
The push to actually deploy these systems in the real world forced researchers to
build leaner, custom networks. Rahman \emph{et al.} (2019) constructed an end-to-end
custom CNN from scratch~\cite{rahman2019}. Pushing the envelope further, Fuhad
\emph{et al.} (2020)~\cite{fuhad2020} and Masud \emph{et al.} (2020)~\cite{masud2020}
designed ultra-lightweight architectures crafted specifically for smartphone
diagnostics and IoT gadgets.

\textbf{Problem Solved:} The bulky computational requirements of ensemble models that
made field deployment virtually impossible.\\
\textbf{How it was Solved:} The teams engineered ultra-light CNNs meant to look
directly at raw red blood cell patches. They also scrubbed mislabeled anomalies from
the dataset so the algorithms could run rapidly on cheap hardware.\\
\textbf{Limitations:} The major trade-off was a total loss of interpretability. These
tiny models operated entirely as ``black boxes,'' leaving doctors in the dark
regarding whether the AI was actually seeing a \emph{Plasmodium} parasite or merely
getting tricked by a random speck of staining dye.

\section{Explainable AI (XAI) in Parasitology}
To win back clinical trust and open up the black box, Rajaraman \emph{et al.} (2018)
integrated Explainable AI (XAI) techniques directly into the diagnostic
pipeline~\cite{rajaraman2018xai}.

\textbf{Problem Solved:} Doctors inherently distrusted lightweight CNNs because they
couldn't see how the AI was making its decisions.\\
\textbf{How it was Solved:} By measuring occlusion sensitivity and generating Class
Activation Maps (CAM), researchers could literally look at what the deepest
convolutional filters were focusing on---proving the model was tracking the actual
parasite rather than background noise.\\
\textbf{Limitations:} Even with visual heatmaps, standard CNNs harbored a built-in
structural flaw. Max-pooling operations aggressively shrink feature maps, which
completely shreds the actual physical shape, spatial hierarchy, and relative
orientation of the parasite.

\section{Advanced Spatial and Contextual Modeling}
To prevent CNN pooling layers from destroying crucial spatial data, Madhu \emph{et
al.} (2023) introduced an imperative Inception-based Capsule Neural Network
(CapsNet)~\cite{madhu2023}. Around the same time, Masud \emph{et al.} (2022) applied a
Multiheaded attention-based Compact Convolutional Transformer (CCT)~\cite{masud2022}.

\textbf{Problem Solved:} The destruction of a parasite's geometric structure and
spatial layout caused by traditional max-pooling.\\
\textbf{How it was Solved:} CapsNets used dynamic routing mechanisms to output vectors
that kept the exact spatial pose intact. Meanwhile, the Transformers utilized
self-attention mechanisms to evaluate the global context of the entire cell all at
once.\\
\textbf{Limitations:} While these models pushed classification accuracy to an
incredible 99.25\%, the complex math involved was brutally heavy. This effectively
undid all previous progress toward edge deployability. Furthermore, the networks were
still stuck classifying perfectly pre-cropped, single cells.

\section{Multi-Cell Object Detection}
Out in real-world clinical settings, a microscope slide isn't a neat grid of isolated
cells; it is a chaotic, overlapping mess. To bridge this practical gap, Zedda \emph{et
al.} (2023) adapted the YOLOv8 object-detection architecture~\cite{zedda2023}.

\textbf{Problem Solved:} Real clinics do not have the time to perfectly crop single
cells before analyzing them.\\
\textbf{How it was Solved:} They removed the prediction heads meant for massive
objects and swapped in a Parasite Attention Module (PAM). This grabbed high-resolution
details from earlier backbone layers to successfully spot tiny ring-stage parasites
across a wide field of view.\\
\textbf{Limitations:} A persistent hurdle was the model's fragility when faced with
extreme data scarcity. It was also highly sensitive to certain data augmentations; for
instance, applying a shear transform would warp the tiny parasites so much that the
model simply failed to recognize them.

\section{Hybrid Architectures and Synthetic Data Generation}\label{sec:synthetic}
Trying to solve the fundamental data shortage without making the model too slow for
rural clinics, Mmileng \emph{et al.} (2025) leaned on advanced spatial augmentations to
generate a massive synthetic dataset of over 606,000 images~\cite{mmileng2025}.

\textbf{Problem Solved:} The constant tug-of-war between needing massive datasets for
advanced models and requiring extreme computational efficiency for resource-limited
hardware.\\
\textbf{How it was Solved:} The researchers trained a ConvNeXt V2 Tiny model---which
blends Vision Transformer hierarchies with CNN speed---on this giant synthetic
dataset, polishing it off with label smoothing and an AdamW optimizer.\\
\textbf{Limitations:} While ConvNeXt represents the current state-of-the-art, relying
so heavily on fake training data and pre-trained weights is risky. It demands rigorous
ongoing validation to ensure the model doesn't suffer a domain shift when exposed to
unfamiliar geographical staining variations or novel parasite strains in the real
world.

\section{Lightweight Attention Networks}\label{sec:danet}
Rather than scaling models up, Ali \emph{et al.} (2025) asked how much attention a
network actually needs, and built DANet, a Dilated Attention Network with only about
2.3 million parameters~\cite{ali2025danet}. Each of its nine Dilated Attention Blocks
stacks $3\times3$ convolutions with dilation rates of 1, 2 and 3, pools the result in
two ways (average and max), and turns it into an attention map that re-weights the
block's input.

\textbf{Problem Solved:} The tug-of-war between the accuracy of heavy CNN--transformer
hybrids and the tiny compute budget of point-of-care hardware, made worse by
low-contrast smears in which parasites have blurry, indistinct borders.\\
\textbf{How it was Solved:} Dilated convolutions widen the receptive field without
adding parameters, while the attention map highlights the parasite region inside
low-contrast cells. DANet reports 97.95\% accuracy and a 97.86\% F1-score on the NIH
cells, runs on a Raspberry Pi~4, and releases its code publicly.\\
\textbf{Limitations:} The evaluation still rests on a random, cell-level split of a
single malaria dataset, so the same patients appear in both the training and the test
sets. Interpretability is shown only through qualitative Grad-CAM examples, and the
network is trained without any regularization on the strength of a no-overfitting
claim.

\section{Research Gaps and Opportunities}\label{sec:gaps}
Even with sky-high accuracy in controlled laboratory tests, we aren't quite ready to
put these models in real clinics without supervision. Phase 1 research flags a few
major roadblocks:
\begin{itemize}
\item \textbf{Data Leakage via Cell-Level Splitting:} A surprisingly common
methodological flaw is randomly splitting the NIH dataset at the individual cell
level, rather than grouping it by patient. When models train this way, they often just
memorize patient-specific staining artifacts instead of learning generalized parasite
morphology. This creates a dangerous illusion of high accuracy.
\item \textbf{Domain Generalization:} If you train a model strictly on the NIH images
or their synthetic clones, it tends to suffer catastrophic performance drops the
second you show it a blood smear from an unrepresented region with slightly different
dye properties or novel parasite strains.
\item \textbf{Transition to Instance-Level Detection:} For true clinical utility, the
field has to move beyond analyzing perfectly isolated cells. The AI needs to scan
entire, messy thick and thin blood smears to detect multiple instances of parasites in
real-time.
\item \textbf{Clinically Validated XAI:} Sure, basic Grad-CAM heatmaps look nice, but
they aren't enough. We need rigorous, medically validated explainability that proves
the network's spatial attention correlates directly with actual parasitological
markers---like hemozoin pigment or chromatin dots---before clinicians will fully trust
the output.
\end{itemize}

\begin{thebibliography}{99}
\bibitem{who2023} World Health Organization, ``World malaria report 2023,'' Geneva,
Switzerland, 2023.
\bibitem{rajaraman2018} S. Rajaraman \emph{et al.}, ``Pre-trained convolutional
neural networks as feature extractors toward improved malaria parasite detection in
thin blood smear images,'' \emph{PeerJ}, vol.~6, p.~e4568, 2018.
\bibitem{rajaraman2019} S. Rajaraman, S. Jaeger, and S. K. Antani, ``Performance
evaluation of deep neural ensembles toward malaria parasite detection in thin-blood
smear images,'' \emph{PeerJ}, vol.~7, p.~e6977, 2019.
\bibitem{madhu2023} G. Madhu \emph{et al.}, ``Intelligent diagnostic model for malaria
parasite detection and classification using imperative inception-based capsule neural
networks,'' \emph{Scientific Reports}, vol.~13, no.~1, p.~13531, 2023.
\bibitem{masud2022} M. Masud \emph{et al.}, ``Explainable Transformer-Based Deep
Learning Model for the Detection of Malaria Parasites from Blood Cell Images,''
\emph{Sensors}, vol.~22, no.~12, p.~4373, 2022.
\bibitem{rahman2019} A. Rahman \emph{et al.}, ``Improving Malaria Parasite Detection
from Red Blood Cell using Deep Convolutional Neural Networks,'' \emph{arXiv preprint}
arXiv:1907.10418, 2019.
\bibitem{fuhad2020} K. M. F. Fuhad \emph{et al.}, ``Deep Learning Based Automatic
Malaria Parasite Detection from Blood Smear and its Smartphone Based Application,''
\emph{Diagnostics}, vol.~10, no.~5, p.~329, 2020.
\bibitem{masud2020} M. Masud \emph{et al.}, ``Leveraging Deep Learning Techniques for
Malaria Parasite Detection Using Mobile Application,'' \emph{Wireless Communications
and Mobile Computing}, vol.~2020, pp.~1--15, 2020.
\bibitem{ali2025danet} A. Ali, R. Pal, I. Dey, E. Cuevas, M. Perez-Cisneros, and
R. Sarkar, ``DANet: a lightweight dilated attention network for malaria parasite
detection,'' \emph{Scientific Reports}, vol.~15, p.~36689, 2025.
\bibitem{tellez2019} D. Tellez, G. Litjens, P. B\'andi, W. Bulten, J.-M. Bokhorst,
F. Ciompi, and J. van der Laak, ``Quantifying the effects of data augmentation and
stain color normalization in convolutional neural networks for computational
pathology,'' \emph{Medical Image Analysis}, vol.~58, p.~101544, 2019.
\bibitem{rajaraman2018xai} S. Rajaraman \emph{et al.}, ``Visualizing and Explaining
Deep Learning Models for Malaria Parasite Detection,'' \emph{Medical Image
Understanding and Analysis}, 2018.
\bibitem{zedda2023} L. Zedda, A. Loddo, and C. Di Ruberto, ``YOLO-PAM:
Parasite-Attention-Based Model for Efficient Malaria Detection,'' \emph{Journal of
Imaging}, vol.~9, no.~12, p.~266, 2023.
\bibitem{mmileng2025} T. Mmileng \emph{et al.}, ``Synthetic Data Generation and Hybrid
Architecture Modeling for Edge-Deployed Malaria Diagnostics using ConvNeXt,''
\emph{IEEE Access}, vol.~13, pp.~100--115, 2025.
\end{thebibliography}

\end{document}
